% Version 0.26.  Complete weighted intersection suppliers and the paired-count cut.
\section{Weighted intersection lengths and the symplectic paired-count frontier}
\label{sec:symplectic-count-closure}

\status{SP1 proves the literal \texttt{LemmaM'}.  SP2--SP6 prove the
literal \texttt{LemmaM}, including its finite orbit cover and its factor
$2$, by a three-variable Rees calculation.  These are complete algebra
suppliers for both the rank-three and rank-four routes.  The later
statements separate the finite trace-product arithmetic and coefficient
return from the still unproved construction and identification of the
symplectic deformation model.  Neither generic automorphy nor a package
named \texttt{HypothesesH'} is accepted as a classical citation leaf.}

\subsection{The cotangent lattice and the additive-injection lemma}

Let $O$ be a discrete valuation ring, with uniformizer $\pi$, and let
$x:A\to O$ be an $O$-algebra augmentation.  Put
\[
 \mathfrak a=\ker x,\qquad
 \mathfrak a^{(2)}=\{f\in A:\ cf\in\mathfrak a^2
                         \text{ for some }0\ne c\in O\}.
 \tag{SP.1}
\]
For a grading $A=\bigoplus_{j\in X}A_j$, define
\[
 C_j=\operatorname{im}(A_j\cap\mathfrak a\longrightarrow
                          A/\mathfrak a^{(2)}),\qquad
 C_j(F)=C_j/\langle\overline F_i\rangle_O
 \quad(F_i\in A_j\cap\mathfrak a).
 \tag{SP.2}
\]
Thus $C_j$ and $\ell_O C_j(F)$ are exactly \texttt{cotLat} and
\texttt{cotLen}.  Length is allowed to be $\infty$.
The source's \texttt{SpansCot} means that the $F_i$ span $C_j$ after
inverting $\pi$; it does not mean that they span the integral lattice.

\begin{proposition}[SP1: the literal \texttt{LemmaM'}]
\label{claim:SP1}
Suppose the residue field of $O$ is finite, $X$ is any abelian grading
group, $x(A_j)=0$, and $0\ne c\in O$.  Let $M$ be an $O$-module and
let
\[
 D:M\longrightarrow\operatorname{Hom}_O(A,O/(c))
 \tag{SP.3}
\]
be an injective additive map.  Suppose $D\mu(F_i)=0$ for every $\mu,i$,
and every $D\mu$ is the reduction of an $O$-valued derivation at $x$
which vanishes on all weights other than $j$.  Then
\[
                    \ell_O M\le\ell_O C_j(F).
 \tag{SP.4}
\]
No $O$-linearity of the assignment $\mu\mapsto D\mu$ is assumed.
\end{proposition}

\begin{proof}
An $O$-valued derivation $\delta$ at $x$ kills $\mathfrak a^2$ by the
Leibniz rule.  If $cf\in\mathfrak a^2$ and $c\ne0$, then
$c\delta(f)=0$ in the domain $O$, so it also kills
$\mathfrak a^{(2)}$.  Restricting the reductions to $A_j$ therefore
gives an additive map
\[
 M\longrightarrow\operatorname{Hom}_O(C_j(F),O/(c)).
 \tag{SP.5}
\]
It is injective: equality on $A_j$ implies equality on every homogeneous
component, because both maps vanish on the other components; every
element of $A$ is a finite sum of homogeneous elements.  Injectivity
of (SP.3) then identifies the two elements of $M$.

If $C_j(F)$ has infinite length, (SP.4) is immediate.  Otherwise the
structure theorem over a DVR gives
$C_j(F)\simeq\bigoplus_iO/(\pi^{a_i})$.  Writing $b=v_O(c)$, the
target of (SP.5) has length $\sum_i\min(a_i,b)\le\sum_i a_i$ and
cardinality $q^{\sum_i\min(a_i,b)}$, where $q=|O/\pi|>1$.
The additive injection makes $M$ finite.  Every finite $O$-module has
cardinality $q^{\ell_O M}$, so comparison of cardinalities proves
(SP.4).  This also explains the finite-residue-field hypothesis in the
source: an additive injection alone is what is available.
\end{proof}

\subsection{The exact weighted intersection problem}

For SP2--SP6, $O$ is a complete DVR of characteristic zero,
$E=\operatorname{Frac}(O)$, and $K$ is an algebraically closed field
with a faithful $O$-algebra structure.  A valuation
$v:K^\times\to\mathbb R$ is normalized by $v(\pi)=1$ and
$v(c)=v_O(c)$ for $c\in O\setminus\{0\}$; these are exactly the
nonzero-input clauses of \texttt{IsNVal}.  Let
$A=\bigoplus_{j\in\mathbb Z}A_j$, $B=A_0$, where $B$ is complete
Noetherian local, $O\to B$ is local, and the induced residue-field map
is an isomorphism.  The last three conditions are \texttt{IsCO O B}.
Assume that $A$ is generated over $B$ by finitely many elements of
weights $-2,-1,1,2$, and that $x(A_j)=0$ for $j\ne0$.
These are precisely the remaining clauses of \texttt{IsWRing}.

Fix integers $r\ge1$, $s,r',s',m\ge0$ and equations
\[
\begin{array}{c|ccccc}
 \text{equations}&F_0,\ldots,F_r&G_1,\ldots,G_s&E_1,\ldots,E_{r'}
       &E'_1,\ldots,E'_{s'}&Z_1,\ldots,Z_m\\
 \text{weights}&-1&-2&1&2&0 .
\end{array}
 \tag{SP.6}
\]
All $x(Z_i)$ vanish, and $I=(F,G,E,E',Z)$.  Write $S(K)$ for the
$O$-algebra maps $\xi:A/I\to K$.  Such a point is \emph{two-sided}
if it is nonzero on some negative and on some positive weight.
The action is $(t\xi)(f_j)=t^j\xi(f_j)$.

We use the following exact five hypotheses, with their source names
recorded to make the endpoint checkable.
\begin{enumerate}
\item[(M1)] There is a weight-preserving formal chart
$\phi:A\to E[[z_i]]$ with constant term $x$, whose linear coordinate
functionals are independently realized by elements of $A$.  Their
common kernel on $\ker x$ is $\mathfrak a^{(2)}$.  The numbers of
coordinates of weights $-1,-2,1,2,0$ are respectively
$r,s,r'+1,s'+1,m$, and there are no other weights.  The equations
$F,G,Z$ span $C_{-1},C_{-2},C_0$ rationally.
\item[(M2)] The image of $B$ in $A/I$ is a finite $O$-module.
\item[(M3)] Over any fixed point $x_1\ne x$ of $S(K)$, the fibre of
restriction to $B$ cannot contain both a nontrivial point supported
only in nonnegative weights and a nontrivial point supported only in
nonpositive weights.
\item[(M4)] At every two-sided point of $S(K)$, any two $K$-valued
point derivations annihilating $I$ are linearly dependent.
\item[(M5)] The chosen $a\in B$ has $x(a)=0$ and $\xi(a)\ne0$ at
every two-sided point of $S(K)$.
\end{enumerate}
The statement called \texttt{LemmaM} is that there is a finite set
$T$ consisting entirely of two-sided points of $S(K)$ and covering
all their orbits, that
$\ell_OC_{-1}(F)<\infty$, and, for \emph{every} such cover,
\[
                    \ell_OC_{-1}(F)\le2\sum_{\xi\in T}v(\xi(a)).
 \tag{SP.7}
\]
A cover may contain more than one representative of an orbit.

We record two elementary finiteness observations used below.  Each
$A_j$ is a finite $B$-module: in a monomial in the chosen generators,
remove submonomials of total weight zero; a remaining monomial of
fixed weight has bounded length, since the only nonzero generator
weights are $\pm1,\pm2$.  The same grouping, separating a pair of
weight-one factors or one weight-two factor, gives
\[
                       A_{2n}=A_2^n,\qquad A_{-2n}=A_{-2}^n
                       \quad(n\ge0),
 \tag{SP.8}
\]
where products are products of $B$-submodules of $A$.
For example, a mixed monomial is reduced by removing opposite pairs
or triples $(2,-1,-1)$ and $(-2,1,1)$; an even one-sided remainder
partitions into blocks of weight $2$ or $-2$.
Also $B$ acts on every $C_j$ through $x:B\to O$, since
$(b-x(b))f\in\mathfrak a^2$ for $f\in\ker x$.  Thus $C_{-1}$ is
finite over $O$, torsion free by (SP.1), and has rank $r$ by the
injective and surjective linear-term clauses of the chart.

\begin{lemma}[SP2: the multigraded Euler calculation]
\label{claim:SP2}
Let $R$ be a Noetherian $\mathbb N^k$-graded $B$-algebra with
$R_0=B$, generated by finite $B$-modules in degrees
$\mathbf e_1,\ldots,\mathbf e_k$, and put
$R_{++}=\prod_i(R_{\mathbf e_i})R$.  Let $M$ be a finite graded
$R$-module, and let $\mathbf t=(t_1,\ldots,t_h)$ be homogeneous
sections of nonnegative multidegrees.  Suppose
\[
        (\mathfrak m_BR_{++})^N
                  \subset\operatorname{Ann}_R(M)+(\mathbf t)
                  \quad\text{for some }N.
 \tag{SP.9}
\]
For all sufficiently large multidegrees $d$, the Koszul homology
pieces have finite $B$-length.  Their alternating sum
$\chi(\mathbf t;M)_d$ is additive in graded exact sequences,
unchanged by permutation of the sections, and nonnegative.  It is
monotone under a graded quotient, and
\[
 \chi(\mathbf t,\mathbf u;M)_d
                 \le\chi(\mathbf u;M/(\mathbf t)M)_d.
 \tag{SP.10}
\]
Suppose in addition that $\pi$ is injective on $M$, that $B$ acts on
$M/\pi M$ through its residue field, and that $y_1,\ldots,y_n$ all
have degree $\alpha$.  If the section list
$(\pi^{b_i}y_i)_i,\mathbf u$ satisfies (SP.9), then
\[
 \chi((\pi^{b_i}y_i)_i,\mathbf u;M)_d
 \ \ge\ \Big(\sum_i b_i\Big)
       \Delta_\alpha^{n-1}\prod_{u\in\mathbf u}\Delta_{\deg u}
                  \ell_B(M/\pi M)_d .
 \tag{SP.11}
\]
Here $\Delta_\beta h(d)=h(d)-h(d-\beta)$ and every assertion is
eventual in \emph{all} coordinates of $d$.
\end{lemma}

\begin{proof}
Every Koszul homology module is finite graded and is annihilated by
$\operatorname{Ann}(M)+(\mathbf t)$.  If its homogeneous generators
have degrees $d_i$, then, once every coordinate of $d$ exceeds those
of the $d_i$ by $N$, its degree-$d$ part is generated using
$R_{++}^N$.  Consequently $\mathfrak m_B^N$ kills that part.  It is
finite over $B$, hence has finite length.  Exact-sequence additivity
follows from the homology long exact sequence.  Permuting a Koszul
basis proves invariance under permutations.

For completeness, positivity here does not invoke an unspecified
intersection-multiplicity theorem.  A finite graded module has a
finite filtration with factors $(R/\mathfrak p)(-b)$ for homogeneous
prime ideals $\mathfrak p$.  Indeed, choose a nonzero homogeneous
element with maximal annihilator; its annihilator is homogeneous
prime, and iterate in the quotient.  Homogeneous primality implies
primality by comparing highest terms in a translation-invariant
ordering of $\mathbb Z^k$; Noetherianity makes the iteration finite.
Additivity reduces positivity to a factor $R/\mathfrak p$.
If $R_{++}\subset\mathfrak p$, its large graded pieces vanish.
Otherwise there is a section outside $\mathfrak p$ of every
nonnegative multidegree: multiply suitable sections of degrees
$\mathbf e_i$ outside $\mathfrak p$.

Induct on the number of sections.  If the first section $t$ is
nonzero in $R/\mathfrak p$, it is a nonzerodivisor, and its Koszul
complex reduces to the quotient by $t$; apply the induction
hypothesis to the remaining list.  If $t=0$ there, its Euler
contribution is the difference between the Euler functions of
$R/\mathfrak p$ and its shift by $\deg t$.  Choose a section $s$ of
that degree outside $\mathfrak p$.  The exact sequence for
multiplication by $s$ identifies this difference with the Euler
function on $(R/\mathfrak p)/s$.  The remaining list is admissible
because $t$ was already zero.  Induction again gives nonnegativity.
The case of an empty list is simply a length.  This proves
positivity without an infinite-residue-field assumption.

Monotonicity follows from additivity and positivity on the kernel.
The recursion
\[
 \chi(t,\mathbf u;M)_d=
 \chi(\mathbf u;M/tM)_d-
 \chi(\mathbf u;M[t](-\deg t))_d
 \tag{SP.12}
\]
and positivity give (SP.10), one section at a time.

Finally the kernel--cokernel exact sequence for a product of two
endomorphisms, or the equivalent two mapping cones, gives
\[
 \chi(\pi^b y,\mathbf u;M)
       =b\,\chi(\pi,\mathbf u;M)+\chi(y,\mathbf u;M).
 \tag{SP.13}
\]
There is no grading shift in the first term because $\pi$ has degree
zero.  Injectivity of $\pi$ identifies that term with
$\chi(\mathbf u;M/\pi M)$; its graded terms have finite length, so
Euler--Poincar\'e identifies it with the indicated iterated difference
of their lengths.  Peel the factors $\pi^{b_i}$ successively,
permuting each $y_i$ to the end of the list after it is peeled.  The
degrees of the other $n-1$ sections remain $\alpha$.  Each contributes
$b_i$ times the same difference in (SP.11); the residual Euler
function is nonnegative.  Replacing a section by its divisor only
enlarges the section ideal, so admissibility holds at every step.
\end{proof}

\subsection{The Rees support argument and the two-sided fibres}

By Smith normal form, change the $F_i$ by an invertible $O$-linear
change of their $r+1$ generators, without changing their ideal or
their cotangent span, so that
\[
 F_i\equiv\pi^{c_i}y_i\pmod{\mathfrak a^{(2)}}\quad(0\le i<r),
 \qquad H:=F_r\in\mathfrak a^{(2)},\qquad
 \ell_OC_{-1}(F)=\sum_{i<r}c_i .
 \tag{SP.14}
\]
Here the classes of $y_i\in A_{-1}$ form an $O$-basis of $C_{-1}$.
The extra generator is obtained from the rank-one kernel of the map
$O^{r+1}\to\langle\overline F_i\rangle_O$; the target is free and
the map splits.  In particular this change really preserves the
integral span of the equations.

Introduce three finite $B$-submodules of $A$ and their multigraded
Rees algebra:
\[
 \begin{gathered}
 N_0=A_{-2},\quad N_1=A_2,\quad
 N_2=BH^2+A_{-2}^2A_2\subset A_{-2},\\
 R=\bigoplus_{d\in\mathbb N^3}L_dT^d\subset A[T_0,T_1,T_2],
 \qquad L_d=N_0^{d_0}N_1^{d_1}N_2^{d_2}.
 \end{gathered}
 \tag{SP.15}
\]
The section list is
\[
 \mathbf t=(F_i^2T_0)_{i<r},\ (G_iT_0)_i,\
                  (E_i^2T_1)_i,\ (E'_iT_1)_i,\ H^2T_2.
 \tag{SP.16}
\]
Put $D_0=(A_{-2}A_2)B\subset B$ and $I_0=I\cap B$.

\begin{lemma}[SP3: support and admissibility on this Rees algebra]
\label{claim:SP3}
Assume (M2), (M3), and rational spanning by $F$ and $G$.
Let $J'$ be any homogeneous ideal of $R$ containing $\mathbf t$
and the $Z_i$.  For all sufficiently large $d$,
$M_d=(R/J')_d$ is a finite $O$-module.  Every nonzero generic
localization of $M_d$ is at a prime
$\mathfrak p=\ker(\sigma:B\to K)$ with
\[
                     I_0\subset\mathfrak p,
                     \qquad D_0\not\subset\mathfrak p.
 \tag{SP.17}
\]
Such a $\sigma$ is the restriction of a two-sided point of $S(K)$.
If also (M5) holds, then for every homogeneous ideal $J\subset R$ containing
the $Z_i$, the list $(\mathbf t,a)$ satisfies (SP.9) on $R/J$.
\end{lemma}

\begin{proof}
We give the prime argument, including the use of (M3).  Let $P$ be a
prime of $R$ containing $\mathbf t,Z$ but not $R_{++}$.  Choose
$f\in N_0$, $g\in N_1$, $h\in N_2$ such that
$fT_0,gT_1,hT_2\notin P$.  The relation
\[
             (H^2T_0)(hT_2)=(hT_0)(H^2T_2)
 \tag{SP.18}
\]
shows that $H^2T_0\in P$, so all $F_i^2T_0$ belong to $P$.
Let $u=(fT_0)(gT_1)\notin P$.  If $cF_i$ has weight zero, then
\[
              (cF_i)^2u=(F_i^2T_0)(c^2fgT_1)\in P.
\]
For $cG_i$ use $(cG_i)u=(G_iT_0)(cfgT_1)$; the analogous identities
for $E_i^2T_1$ and $E'_iT_1$ interchange the two variables.  Terms
involving $Z_i$ already lie in $P$.  Taking weight-zero components
in an expression for $b\in I$ now shows $I_0\subset P\cap B$.

By (M2), $B/I_0$ is finite over $O$.  Its primes are either the
closed prime $\mathfrak m_B$ or horizontal primes whose residue
fraction fields are finite over $E$.  In the latter case choose an
embedding into $K$, obtaining $\sigma:B\to K$ with
$\ker\sigma=P\cap B$.  We show that $D_0\subset\ker\sigma$ is
impossible.

Here is an explicit construction of the two forbidden one-sided
points.  For $\epsilon\in\{-1,1\}$ use the Rees direction $T_0$
or $T_1$, respectively, and set
\[
 \iota_n(b)=\operatorname{pr}_{2\epsilon n}(b)T_\epsilon^n,
 \qquad
 J_\epsilon=\{b\in A:\ \iota_n(cb)\in P
                         \text{ for all }c\in A,n\ge0\}.
 \tag{SP.19}
\]
The notation $T_{-1}=T_0$, $T_1=T_1$ is used only in (SP.19).
By (SP.8), every displayed monomial lies in $R$.
The set $J_\epsilon$ is an ideal.  Choose $q\in A_{2\epsilon}$
with $qT_\epsilon\notin P$.  For homogeneous elements the identities
\[
 \iota_{n+n'}(bc)=\iota_n(b)\iota_{n'}(c),\qquad
 \iota_{2n}(b^2)=\iota_n(b)^2
 \tag{SP.20}
\]
give the following facts by primality: $b^2\in J_\epsilon$ implies
$b\in J_\epsilon$; $qb\in J_\epsilon$ implies
$b\in J_\epsilon$; and $bT_\epsilon\in P$ for
$b\in A_{2\epsilon}$ implies $b\in J_\epsilon$.
To verify the last assertion for a test factor of weight
$2\epsilon n-2\epsilon$, multiply it by $q$ and cancel
$qT_\epsilon$ in (SP.20).  Likewise
$P\cap B\subset J_\epsilon$ and
$bq^n\notin J_\epsilon$ if $b\notin P\cap B$.
Arbitrary test factors reduce to homogeneous ones by their finite
weight decomposition.

If $D_0\subset\ker\sigma$, then $qA_{-2\epsilon}$ lies in
$P\cap B$, so cancellation puts $A_{-2\epsilon}$ in
$J_\epsilon$.  Squaring and (SP.8) put every weight of the opposite
sign in that ideal.  The section equations in the chosen direction
put the equations of the same sign in it, using squares for weight
$\epsilon$; (SP.18) includes the last $F$.  Thus
$I\subset J_\epsilon$.
The ideal $J_\epsilon$ is disjoint from the multiplicative set
$\{bq^n:\sigma(b)\ne0\}$.  Put
$L=\operatorname{Frac}(B/\ker\sigma)$.
The algebra $((A/J_\epsilon)\otimes_B L)[1/q]$ is a nonzero
finite-type $L$-algebra: disjointness proves nonvanishing, and the
finite generation of $A$ over $B$ proves finite type.
Faithfully flat extension from $L$ to $K$ preserves both properties.
The Nullstellensatz applied to that algebra produces a point
$\xi_\epsilon$ over $\sigma$
annihilating $I$ and all weights of the opposite sign, with
$\xi_\epsilon(q)\ne0$.

There is also a fixed point over $\sigma$, namely
$\xi_0(b)=\sigma(\operatorname{pr}_0b)$.  It is multiplicative:
the square of a product of opposite nonzero weights is a product of
elements of $A_{-2}A_2$ by (SP.8), so its value under $\sigma$ is
zero.  The equations of nonzero weight and the $Z_i$ vanish at
$\xi_0$.  If $\xi_0\ne x$, the two constructed points contradict
(M3).  If $\xi_0=x$, use the negative point.  The functional
$b\mapsto\xi_{-1}(\operatorname{pr}_{-1}b)$ kills
$\mathfrak a^2$: in a product of two augmentation-zero homogeneous
elements of total weight $-1$, at least one has nonnegative weight
and evaluates to zero.  It kills the saturation as well, since
$O\to K$ is injective.  Rational spanning by $F$ forces its
vanishing on $A_{-1}$.  For total weight $-2$, the only additional
case is a product of two weight-$-1$ elements; it now also vanishes.
Rational spanning by $G$ forces vanishing on $A_{-2}$, and (SP.8)
forces vanishing on all negative weights.  This contradicts the
choice of $\xi_{-1}$.  We have proved that a relevant $P$ containing
$D_0$ must contain $\mathfrak m_B$.

These prime statements imply, by Noetherianity, power containments
\[
 (I_0R\,R_{++})^b\subset J',\qquad
 (\mathfrak m_BR\,R_{++})^{b'}\subset J'+D_0R
 \tag{SP.21}
\]
for some $b,b'$.  In large multidegrees they give
$I_0^bM_d=0$ and $\mathfrak m_B^{b'}M_d\subset D_0M_d$.
The first makes $M_d$ finite over $O$: $B/I_0^b$ is finite over
$O$, by its filtration with finite $B/I_0$-modules.  At a horizontal
prime containing $D_0$, $\pi$ is a unit, so the second containment
and Nakayama give $(M_d)_{\mathfrak p}=0$.  This proves (SP.17).

If (SP.17) holds, the degree-zero part of
$K\otimes_B(A/I)$ is $K$, so this finite-type $K$-algebra is
nonzero.  A $K$-point furnished by the Nullstellensatz restricts to
$\sigma$.  A mixed product in $D_0$ has nonzero value, hence this
point is two-sided.  Finally, a relevant prime containing
$J,\mathbf t,a$ and not $\mathfrak m_B$ would yield such a point
with $\xi(a)=0$, contradicting (M5).  Therefore
$\mathfrak m_BR\,R_{++}$ lies in the radical of $J+(\mathbf t,a)$;
a power lies in that ideal, which is (SP.9).
\end{proof}

\begin{lemma}[SP4: finite orbit cover and generic multiplicity one]
\label{claim:SP4}
Assume (M2) and (M4).  The two-sided points of $S(K)$ have finitely
many $K^\times$-orbits.  If $\sigma:B\to K$ and
$D_0\not\subset\ker\sigma$, then every weight piece of $A/I$,
localized at $\ker\sigma$, has $B_{\ker\sigma}$-length at most one.
Replacing all the $F_i$ and $E_i$ by their squares raises this bound
to at most $2^{r+1+r'}$.
\end{lemma}

\begin{proof}
Choose $f\in A_{-2}$ and $g\in A_2$.  Over $E$, the slice
\[
                   C=(A/I)[1/(\pi g)]/(f-1)
 \tag{SP.22}
\]
is a finite-type algebra.  A $K$-point of it is two-sided.  Its Euler
derivation on $A$ is $\mathcal E(b_j)=j\xi(b_j)$ and has
$\mathcal E(f)=-2$.  Any point derivation of $C$ kills $f$.
By (M4) it is linearly dependent with $\mathcal E$; evaluating that
dependence on $f$ forces the Euler coefficient to be zero and then
forces the other derivation to vanish.  After base change to $K$,
every closed point thus has zero cotangent space.  Nakayama makes
its local ring a field.  Nullstellensatz then shows that this
zero-dimensional reduced finite-type algebra is a finite product
of copies of $K$.  Consequently $C$ is a finite \emph{\'etale}
$E$-algebra.

Finite $B$-generating sets of $A_{-2}$ and $A_2$ supply finitely many
such pairs $(f,g)$.  Every two-sided point has both one of the $f$'s
and one of the $g$'s nonzero, by (SP.8).  A square root in $K$
normalizes $f$ to $1$ by the torus action.  The finite point sets of
the slices (SP.22) therefore give a finite orbit cover.

We also need the scheme multiplicity, not just the number of points.
On $U=(A/I)[1/(\pi fg)]$, $f$ is a unit.  Adjoin $z$ with
$z^2=f$.  There is an explicit isomorphism
\[
 U[z,z^{-1}]/(z^2-f)\ \simeq\ C[t,t^{-1}],
 \qquad b_j\longmapsto\overline b_jt^{-j},\quad z\longmapsto t.
 \tag{SP.23}
\]
Its inverse sends $\overline b_j$ to $b_jz^j$.
The left algebra is faithfully flat over $U$ (indeed finite free of
rank two), and the right algebra is reduced.  Thus $U$ is reduced.
This proves the required local reducedness directly, without
assuming that a finite orbit set was already a reduced scheme.

Let $\mathfrak p=\ker\sigma$.  If the localized quotient is zero
there is nothing to prove.  Otherwise the image of $B$ in $A/I$ is
finite over $O$, so its localization at this horizontal prime is an
Artinian local $E$-algebra.  Choose $fg\notin\mathfrak p$.
Its maximal ideal is nilpotent and acts by nilpotents on the reduced
ring in (SP.23), hence acts by zero.  The localized ring is therefore
an algebra over the residue field $L=\operatorname{Frac}(B/\mathfrak p)$,
and its degree-zero part is $L$.  Its even weight pieces have
dimension at most one, since $f$ is an invertible element of degree
$-2$.  If a piece of odd weight $j$ contains $0\ne h$, reducedness
gives $h^2\ne0$.  The even piece of weight $2j$ is one-dimensional,
so $h^2f^j\in L^\times$.  Thus $h$ is invertible; every other
element of weight $j$ is an $L$-multiple of $h$.  This proves the
length-one assertion for every weight.

For the final assertion filter a quotient by $u^2$ using
$uM/u^2M$.  That subquotient is an image of $M/uM$, with the
appropriate weight shift, and the other subquotient is $M/uM$.
Since the length-one bound holds at every weight, each replacement
of one equation by its square multiplies the bound by at most two.
There are $r+1+r'$ such replacements.
\end{proof}

\subsection{The leading form and the exact factor two}

\begin{lemma}[SP5: the lower Euler bound]
\label{claim:SP5}
Under (M1), make the Smith change (SP.14).  There is a
weight-preserving chart $\phi_0$ with only nonzero-weight variables,
with the same linear coordinates in those weights, and
$\phi_0(Z_i)=0$ for all $i$.  Let $J\subset R$ consist of the Rees
polynomials all of whose coefficients lie in $\ker\phi_0$.
If $x(a)=0$ and $(\mathbf t,a)$ is admissible, then for all
sufficiently large $d$,
\[
 2^{r+r'}\sum_{i<r}c_i
              \le\chi(\mathbf t;(R/J)/a(R/J))_d.
 \tag{SP.24}
\]
\end{lemma}

\begin{proof}
The $m$ equations $Z_i$ have linearly independent linear terms in
the $m$ zero-weight variables, because they rationally span that
cotangent space.  Use those equations as the new zero-weight
coordinates.  A formal substitution with invertible linear matrix
has an inverse: after solving the linear terms, solve its
homogeneous terms one total degree at a time.  Each correction has
higher total degree, so the coefficients stabilize.  Uniqueness
also makes the substitution weight preserving.  Set these new
zero-weight coordinates to zero.  This constructs $\phi_0$ and
preserves the linear terms of every remaining variable.

For a monomial $z^\nu$ in these variables let $n_-(\nu)$ and
$n_+(\nu)$ denote its negative and positive weighted degrees.
Every coefficient of a Rees element of degree $d$ has weight
$2(d_1-d_0-d_2)$, and its image under $\phi_0$ has
\[
                    n_-\ge 2d_0+4d_2.
 \tag{SP.25}
\]
For $N_0$ and $N_1$ this follows from the weight.  For $N_2$, the
term $A_{-2}^2A_2$ has this lower bound, while $H$ has no constant
or linear term; a weight-$-1$ monomial without a linear term has
$n_-\ge2$, so $H^2$ has $n_-\ge4$.  Products prove (SP.25).

Let $Q\subset R$ be the ideal obtained by requiring strict inequality
in (SP.25) in every Rees degree.  Leading coefficients show directly
that it is a homogeneous ideal, $J\subset Q$, $a\in Q$, and
$B$ acts on $R/Q$ through $x$.  It is $O$-torsion free: multiplying
a leading coefficient in $E$ by a nonzero scalar cannot kill it.
Moreover (SP.14) gives, in $R/Q$,
\[
                      F_i^2T_0=\pi^{2c_i}y_i^2T_0.
 \tag{SP.26}
\]
Indeed the difference $F_i-\pi^{c_i}y_i$ has no linear term and
has weight $-1$, so every term of the difference of squares has
negative weighted degree strictly greater than $2$.

The degree-$d$ part of $R/Q$ is a finite torsion-free $O$-module.
Its rank is
\[
 H(d)=\#\{\nu:\ n_-(\nu)=2d_0+4d_2,
                       \ n_+(\nu)=2d_1+2d_2\}.
 \tag{SP.27}
\]
To check equality of the rank, rather than only an upper bound,
use the independently realized linear coordinates in the chart.
The leading weight-$\pm1$ forms are independent linear forms; the
leading weight-$\pm2$ forms are independent linear forms in the
weight-two variables plus quadratics in the weight-one variables.
Subtract those quadratics over $E$.  These leading forms therefore
generate the entire polynomial algebra on each sign of variables.
Every monomial of even weighted degree is a product of
weighted-degree-two monomials.  Finally
\[
             A_{-2}^{d_0+2d_2}A_2^{d_1+d_2}\subset L_d
\]
realizes every monomial counted in (SP.27).  This proves full rank
after tensoring with $E$.  No claim that the integral lattice is
unimodular is being used.  Reduction of a finite free $O$-lattice
modulo $\pi$ has length its rank, so
$\ell_B((R/Q)/\pi)_d=H(d)$.

There is a graded surjection $(R/J)/a\to R/Q$.  Monotonicity in
SP2 and the scaling inequality (SP.11), applied to the $r$ sections
in (SP.26), give the lower bound
\[
 2\sum_i c_i\,
       \Delta_0^{r-1+s}\Delta_1^{r'+s'}\Delta_2H(d).
 \tag{SP.28}
\]
All needed admissibility conditions descend under the surjection.
The remaining numerical difference has a short exact calculation.
The generating function for monomials of bidegree $(a,b)$ is
\[
 \frac{1}{(1-X)^r(1-X^2)^s(1-Y)^{r'+1}(1-Y^2)^{s'+1}}.
\]
The three differences in (SP.28) are shifts by $(2,0),(0,2),(4,2)$.
After multiplying by their factors the generating function becomes
\[
 (1+X)^{r-1}(1+Y)^{r'}
               \frac{1-X^4Y^2}{(1-X)(1-Y)(1-Y^2)}.
 \tag{SP.29}
\]
For sufficiently large $a,b$, the coefficient in the last quotient
is
$(\lfloor b/2\rfloor+1)-(\lfloor(b-2)/2\rfloor+1)=1$.
The polynomial factor therefore makes the coefficient exactly
$2^{r-1+r'}$.  Substitution in (SP.28) proves (SP.24).
\end{proof}

\begin{theorem}[SP6: the literal \texttt{LemmaM}]
\label{claim:SP6}
Under exactly (M1)--(M5) and the ring and valuation hypotheses
above, there is a finite cover of the two-sided orbits,
$\ell_OC_{-1}(F)$ is finite, and (SP.7) holds for every finite orbit
cover.  Thus \texttt{LemmaM} has no remaining bespoke input.
\end{theorem}

\begin{proof}
The cotangent finiteness and (SP.14) already prove the finiteness of
the left side, and SP4 supplies a finite cover.  We first arrange
$\phi_0(a)\ne0$.  If it is zero, rational spanning by $F$ and the
nonzero weight-$-1$ linear coordinate imply that some $F_i$ has
nonzero chart image.  Choose $g\in A_1$ with nonzero chart image;
there is such a linear coordinate because its dimension is $r'+1$.
Replace $a$ by $a+gF_i$.  It still has weight zero and augmentation
zero, it agrees with $a$ on $S(K)$, and its chart image is a product
of two nonzero power series in a domain.  This replacement preserves
(M5) and the right side of (SP.7).

For the coefficientwise kernel $J$ of SP5, multiplication by $a$
is injective on $R/J$: its image is multiplication by the nonzero
series $\phi_0(a)$ on each coefficient.  SP3 supplies admissibility.
Apply SP2, first moving $a$ past the other sections, then passing to
their quotient.  Together with SP5 this gives
\[
 2^{r+r'}\sum_i c_i
 \le \chi(\mathbf t;(R/J)/a)_d
 =\chi(\mathbf t,a;R/J)_d
 \le\chi(a;R/(J+(\mathbf t)))_d .
 \tag{SP.30}
\]
We bound the final term on the actual $B$-module
$M_d=(R/(J+(\mathbf t)))_d$.

By SP3, $M_d$ is finite over $O$, and its generic support satisfies
(SP.17).  At such a prime choose $f\in A_{-2}$ and $g\in A_2$
with $fg\notin\mathfrak p$.  They become invertible after
localization, so $N_0,N_1,N_2$ all become free rank-one modules in
their indicated weights; $N_2=N_0$ there since
$f^2g=(fg)f\in A_{-2}^2A_2$.  In degree $d$ with $d_i\ge1$, the
section relations (SP.16) impose every homogeneous multiple of the
equations
\[
                    I^{[2]}=(F_i^2,G_i,E_i^2,E'_i,Z_i).
\]
Indeed one may take that multiple in degree $d-\mathbf e_i$ in
the corresponding Rees direction.  Consequently the localized
weight-$2(d_1-d_0-d_2)$ piece of $A/I^{[2]}$ surjects onto
$(M_d)_{\mathfrak p}$.  SP4 gives
\[
                    \ell_{B_{\mathfrak p}}(M_d)_{\mathfrak p}
                         \le2^{r+1+r'}.
 \tag{SP.31}
\]

Here is the length--valuation formula needed to sum these
multiplicities.  For any finitely generated $O$-module $M$ with a finite
$B$-action and any $b\in B$ invertible on $E\otimes_OM$,
\[
 \ell_B(M/bM)-\ell_B(M[b])
       =\sum_{\sigma}m_\sigma v(\sigma(b)),\qquad
 m_\sigma=\ell_{B_{\ker\sigma}}M_{\ker\sigma}.
 \tag{SP.32}
\]
The sum is over all embeddings of the generic support fields into
$K$.  To prove it, remove the $O$-torsion of $M$; its kernel and
cokernel under any endomorphism have the same finite length, so
contribute zero.  On the remaining free $O$-lattice the left side is
$v_O(\det(b\mid E\otimes_OM))$.  Decomposition of the finite
commutative $E$-algebra acting on this vector space into Artinian
local factors expresses this determinant as a product of field
norms, each raised to the length of the corresponding localized
module.  In characteristic zero each norm is the product over its
embeddings into $K$; applying $v$ gives (SP.32).  Finite $B$-length
equals finite $O$-length because their residue fields agree.

Every embedding in (SP.32) is the restriction of a two-sided point
by SP3, and hence of some representative in the given cover $T$.
By (M5), $a$ is nonzero at all such embeddings and is invertible on
$E\otimes_OM_d$.  Since the image of $B$ on $A/I$ is finite over
$O$, each $\xi(a)$ is integral over $O$.  A root of a monic
polynomial with coefficients of nonnegative valuation cannot have
negative valuation: its leading power would be the unique term
of smallest valuation.  Thus all $v(\xi(a))$ are nonnegative.
Repeated restrictions in $T$ can therefore only increase their
sum.  Equations (SP.31)--(SP.32) imply
\[
            \chi(a;M_d)\le
                    2^{r+1+r'}\sum_{\xi\in T}v(\xi(a)).
 \tag{SP.33}
\]
Combine (SP.30) and (SP.33), divide by $2^{r+r'}$, and use (SP.14).
This is exactly (SP.7), including the source's constant $2$.
\end{proof}

The classical background used in SP1--SP6 is the DVR structure
theorem, finite-dimensional linear algebra, the Nullstellensatz and
Nakayama's lemma.  The special multigraded positivity, prime support,
formal substitution, orbit reducedness and numerical difference
arguments have been proved above.  For the two algebraic-geometric
background statements, see the Stacks Project, Tags
\href{https://stacks.math.columbia.edu/tag/00FV}{00FV} and
\href{https://stacks.math.columbia.edu/tag/00DV}{00DV}; the finite
\'etale conclusion in SP4 is also Tag
\href{https://stacks.math.columbia.edu/tag/00U3}{00U3}.

\subsection{From actual rank-four cocycles to this cotangent lattice}

We now return to the source's original $\Omega$, odd prime $\ell$,
coefficient ring $O=c.O\subset\overline{\mathbb Q}_\ell$, and datum
$D:\texttt{Raw}\ \Omega\ c$.  The statement $D.\texttt{Good}$
includes $e=D.e\ge2$, an open-kernel finite-order character $D.\psi$,
the specified places $D.\mathrm{pl}(p)$, and
$\Psi=D.\psi\chi_\ell^e$.  It also includes
$\overline\Psi(g_0)\ne1$ and a local element at $\ell$ on which
$\chi_\ell$ has infinite order and $D.\psi$ is trivial.
Put
\[
 S_0=D.S,\quad
 T=\{\ell\}\cup\{p\in S_0:p\ne\ell,
                      \ \ell\mid|\psi(I_p)|\},\quad
 W=\operatorname{Frac}(O)/O(\Psi^{-1}).
 \tag{SP.34}
\]
The underlying torsion module in (SP.34) is the source's
\texttt{QO}; its twist, valuations and decomposition subgroups are
unchanged.  Let $\mathcal S$ be the strict Selmer module imposing
trivial restriction to every inertia group and to the decomposition
groups at primes in $S_0$.  Then
$D.\mathrm{sha}=\ell_O\mathcal S$ by definition.

\begin{proposition}[SP7: the actual \texttt{shaCoc} and \texttt{mPrime} comparison]
\label{claim:SP7}
Assume $D.\texttt{Good}$ and RV3.  Then $\mathcal S$ has finite
length.  There is a single $a\in O\setminus\{0\}$ and an additive
injection
\[
          Z:\mathcal S\lhook\joinrel\longrightarrow
                  \{z:G_{\mathbb Q}\to O/(a)\}
 \tag{SP.35}
\]
whose values are exactly of the source's type $D.\texttt{ShaCoc}$:
locally constant $\Psi^{-1}$-cocycles, zero at $g_0$, zero on the
residual-trivial inertia tested by $D.\texttt{UnrT}$, and
coboundaries over $O/(a)$ on every specified decomposition group
in $T$.
For every actual $\mathcal M:\texttt{Model}\ D$, equipped with the
source \texttt{GradedAlgebra} structure on its first-weight-zero
algebra $A$ and its pieces \texttt{Model.gr}, one has
\[
       D.\mathrm{sha}\le
       \ell_OC_{-1}(\mathcal M.F).
 \tag{SP.36}
\]
Thus this particular incoming edge to the trace inequality is closed
once an actual model is supplied; no additional cocycle comparison
or choice-linearity hypothesis remains.
\end{proposition}

\begin{proof}
The strict Selmer module for $S_0$ is a submodule of the one imposing
only the decomposition condition at $\ell$, because $\ell\in S_0$.
RV3 therefore gives finite length.  Choose $0\ne b\in O$ killing
$\mathcal S$ and write $\theta=\Psi^{-1}$.
The element $u=\theta(g_0)-1$ is a unit, since its reduction is
nonzero.  Every cohomology class has the unique normalized cocycle
\[
                 N[z](g)=z(g)-(\theta(g)-1)u^{-1}z(g_0).
 \tag{SP.37}
\]
Changing $z$ by a coboundary leaves (SP.37) unchanged; it vanishes
identically exactly when the class is zero.  Hence $N$ is an
$O$-linear injection into locally constant cocycles and every
$N[z](g)$ is killed by $b$.

For each $v\in T$, choose $s_v\in D_v$ with
$\theta(s_v)\ne1$.  At $\ell$ use the infinite-order cyclotomic
element specified by $D.\texttt{Good}$.  At $v\ne\ell$ the
definition of $T$ makes $\psi(I_v)$ a nontrivial group of order
divisible by $\ell$, while $\chi_\ell$ is unramified there; an
element on which $\psi$ is nontrivial suffices.  Set
\[
                       a=b\prod_{v\in T}(\theta(s_v)-1)\ne0.
 \tag{SP.38}
\]
The normalized cocycle is a local coboundary
$N[z](g)=(\theta(g)-1)m_v$ on $D_v$.  Since
$b(\theta(s_v)-1)m_v=bN[z](s_v)=0$, (SP.38) kills $m_v$ as well
as all the values of $N[z]$.
The map $O/(a)\to\operatorname{Frac}(O)/O$ sending
$\overline y$ to $y/a+O$ identifies its source with the $a$-torsion.
Transport both the cocycle and these local parameters through this
map.  This constructs the additive injection (SP.35), with all its
local and normalization properties.

We check the potentially stronger-looking $\texttt{UnrT}$ clause.
Outside $S_0$ the finite character is unramified; at a prime of
$S_0\setminus T$ its inertia image has order prime to $\ell$.
The cyclotomic character is unramified away from $\ell$.
If the residual diagonal representation is trivial on an inertia
element, then $\theta$ is $1$ modulo the maximal ideal there.  A
root of unity of order prime to $\ell$ that is $1$ modulo that ideal
equals $1$: divide $X^n-1$ by $X-1$ and evaluate at the root, where
the second factor reduces to the unit $n$.
Thus $\theta(g)=1$ for exactly the inertia elements needed in
$\texttt{UnrT}$, and local coboundary gives $N[z](g)=0$.
Conjugating the chosen places gives the same assertion for every
valuation subring in that definition.  A valuation trivial on
$\mathbb Q$ is trivial on its algebraic closure and has trivial
inertia, so it adds no condition.  This verifies the literal
$D.\texttt{ShaCoc}$ statement.

It remains to use the actual model rather than postulate an
injection into cotangent duals.  Its record supplies a surjection
$q:R\to S$ onto the universal representation ring, a
$\mathbb Z^2$-weighted presentation on $R$, and, for every cocycle
$Z\mu$, an $O$-valued point derivation $\delta_\mu$ of weight
$-w_\Psi=(0,-1)$.  It kills $\ker q$ modulo $a$.  On lifts of the
two relevant universal matrix entries its reductions are
\[
 \delta_\mu(r(g)_{20})=D.\mathrm{th}(g)Z\mu(g),\qquad
 \delta_\mu(r(g)_{31})=\Psi(g)Z\mu(g)
                         \pmod a.
 \tag{SP.39}
\]
All other entries have different weights.  These are precisely the
\texttt{der}, \texttt{hom}, \texttt{at0} and \texttt{gen} fields of
$\texttt{Model}$, including the universal quantification over
possible lifts of an entry.

The choices $\delta_\mu$ need not be additive.  Their reductions
are additive.  First the augmentation descends through $q$:
if $q(f)=0$, choose a lift $b$ of the $(0,0)$ entry at the identity;
the \texttt{at0} field applied to $b+f$ and to $b$ gives
$x(b+f)=x(b)$, hence $x(f)=0$.
Now the derivation
$\delta_{\mu+\nu}-\delta_\mu-\delta_\nu$ descends modulo $a$ to
$S$, and is zero on every universal matrix entry by (SP.39) and
the weight restriction.  Associate to it the square-zero algebra
map $S\to O\oplus O/(a)$ with first coordinate the augmentation.
The model's generation field says that two maps to this complete
local $O$-algebra which agree on the universal representation are
equal.  Compare with the zero derivation to deduce vanishing on
all of $S$ and hence all of $R$.
Restriction to the first-weight-zero subalgebra $A\subset R$
therefore gives an additive map
$\mathcal S\to\operatorname{Hom}_O(A,O/(a))$.
It is injective: the $(2,0)$ entry lies in that subalgebra and
the coefficient $D.\mathrm{th}(g)$ in (SP.39) is a unit, so a zero
derivation forces $Z\mu=0$.
Each equation $\mathcal M.F_i$ lies in $\ker q$ and has weight
$-1$ in $A$.  The hypotheses of SP1 now all hold and yield
(SP.36).  This proves the actual source \texttt{mPrime} conclusion
without assuming that an arbitrary choice of lifts was linear.
\end{proof}

\subsection{Returning from auxiliary coefficients to the original character}

\begin{proposition}[SP8: exact character invariants and coefficient return]
\label{claim:SP8}
Let $f:O\hookrightarrow O'$ be an injective compatible inclusion of
coefficient DVRs in $\overline{\mathbb Q}_\ell$, and let
$\psi' = f\circ\psi$ be a Dirichlet character of the same modulus
$N$.  Then order, unit-character kernel, conductor, the parity
condition $\psi(-1)=(-1)^e$, and all the source's quantities
\[
          \operatorname{comp}_r\psi,\quad\theta_0,\quad\theta_1,
          \chi_P,\quad n_0,\quad c_2,\quad n_{\rm II},\quad n_{\rm III}
 \tag{SP.40}
\]
are preserved or commute with $f$ as appropriate.  In particular
membership in $h.X(c,e)$ and either case II or case III is preserved
on the same host and the same $\Omega$.
Together with RV1 this removes all auxiliary root-of-unity
assumptions from a bound expressed using the normalized Selmer
length: the identical bound holds over the original $O$.
\end{proposition}

\begin{proof}
Injectivity gives $f(\psi(u))=1$ if and only if $\psi(u)=1$,
for every unit $u$ modulo $N$, and $\psi'^n=1$ if and only if
$\psi^n=1$.  Thus the unit kernels and orders agree.
A character factors through $d\mid N$ exactly when its unit kernel
contains the kernel of $(\mathbb Z/N)^\times\to(\mathbb Z/d)^\times$.
Its complete set of possible conductors, and hence its conductor,
is unchanged.  The operations $\operatorname{comp}_r$ are pullbacks
by the fixed primary projection of $(\mathbb Z/N)^\times$; they
commute with any coefficient map.
The prime-to-$\ell$ part $\theta_0=\ell\mathrm{Prime}(\psi)$ is
defined by group powers determined by the order.  It too commutes
with $f$, as do
\[
 \theta_1=\theta_0(\operatorname{comp}_2\theta_0)^{-1},\qquad
 \chi_P=\theta_0(\operatorname{comp}_\ell\theta_0)^{-1}.
\]
The definitions
\[
 \begin{aligned}
 f_0&=\operatorname{ord}_{\operatorname{lcm}(4,\operatorname{ord}\theta_1)}(\ell),
 &n_0&=v_2(\ell^{f_0}-1),\\
 c_2&=v_2(\operatorname{cond}(\operatorname{comp}_2\psi))
 \end{aligned}
 \tag{SP.41}
\]
then prove their preservation.  So do the defining ceiling logarithms
for $n_{\rm II}$ and $n_{\rm III}$, written in (SP.48) below.
Injectivity also gives
$f(\psi(-1))=(-1)^e\Longleftrightarrow\psi(-1)=(-1)^e$.
Host character membership is a unit-kernel condition together with
this parity.  Case II uses $e\ge2$, the parity of
$\operatorname{ord}\chi_P$, and the fixed numerical condition on
$\ell,h,e$; case III uses that same fixed numerical condition,
odd order of $\psi$, and $\psi\ne1$.  All are therefore preserved.

For an individual $\psi$, adjoin a primitive
$2^{n_{\rm II}+n_{\rm III}}$-th root of unity to the finite
coefficient field and take its compositum with the original field.
Its ring of integers $O'$ is again an allowed coefficient ring,
contains the required roots for each separate exponent, and
contains $O$ injectively.  RV1 identifies the original strict
Selmer module after this extension with its scalar extension.
If $e_{O'/O}$ is the relative ramification index, then
\[
 \ell_{O'}\mathcal S'=e_{O'/O}\ell_O\mathcal S,
 \qquad v_{O'}(\ell)=e_{O'/O}v_O(\ell),\qquad
 \frac{\ell_{O'}\mathcal S'}{v_{O'}(\ell)}
       =\frac{\ell_O\mathcal S}{v_O(\ell)}.
 \tag{SP.42}
\]
Finiteness is equivalent in both directions.  Thus a bound proved
over $O'$ gives the same bound over $O$, with its original
character, conductor, parity, case, and normalization.  The chosen
coefficient extension may depend on $\psi$; a constant selected
before $\psi$ is not changed by this construction.
\end{proof}

\subsection{The exact arithmetic remaining after a finite trace catalogue}

The existence of the following catalogue is a retained input, not a
synonym for a general automorphy theorem.  Its quantifier order is
part of the input.  There must be integers $c_*,A_*\ge0$ such that,
for every $e\ge2$, $M\ge1$, and prime $p\ge5$ with $p\nmid M$,
there is a finite set $\mathcal D$ of complex multisets
$\beta_s$ of cardinality four with the following properties:
\begin{enumerate}
\item $|\mathcal D|\le c_*[M(e+2)]^{A_*}$, and every member of every
$\beta_s$ is an algebraic integer.
\item Every automorphism of $\mathbb C$ permutes the multisets
$\beta_s$, with multiplicity, and
\[
 \left|\sum\beta_s\right|
       \le (3\sqrt p+1)p^{(e+1)/2}.
 \tag{SP.43}
\]
\item For every odd prime $\ell\ne p$, every original algebraic
closure $\Omega/\mathbb Q$, every isomorphism
$\iota:\overline{\mathbb Q}_\ell\simeq\mathbb C$, and every
representation satisfying the literal \texttt{AdmRAlt} conditions
at level $M$, weights $(e+1,e,1,0)$ and duality weight $e+1$, there
is a label $s\in\mathcal D$.  Equal labels imply representation
isomorphism.  At every allowed place above $p$ and every arithmetic
Frobenius $\sigma$ there,
\[
                  \iota(\operatorname{tr}r(\sigma))=\sum\beta_s.
 \tag{SP.44}
\]
\end{enumerate}
Here \texttt{AdmRAlt} includes irreducibility, the literal
\texttt{IsAutQA} condition, \texttt{Level}, the stated cyclotomic
dual twist, and \texttt{SafeRA}.  In the source the catalogue is
first expressed using geometric Frobenius; its inverse is used in
(SP.44).  Both this inverse convention and the classification of
all the actual representations are required.  The catalogue is
chosen before $\ell,\Omega,\iota$ and the representations.

\begin{proposition}[SP9: the integral product and its uniform trace bound]
\label{claim:SP9}
The preceding catalogue input supplies the literal
\texttt{NumberSp4Alt}, with
\[
                    C=c_*2^{A_*+1},\qquad A=A_*+1.
 \tag{SP.45}
\]
More precisely, after $e,M,p,k>0$ and three $k$-th roots of unity
$z_1,z_2,z_3$ have been fixed, there is a nonzero integer $N_*$,
chosen before $\ell,\Omega,\iota$ and any finite representation
family, such that
\[
 \begin{aligned}
 \log|N_*|&\le C\varphi(k)[M(e+1)]^A\log p,\\
 |N_*|_\ell&\le
 \prod_{r\in\mathcal X}
 \left|\operatorname{tr}r(\sigma)
       -\iota^{-1}\bigl(p^{e+1}z_1+p^ez_2+pz_3+1\bigr)\right|_\ell.
 \end{aligned}
 \tag{SP.46}
\]
Every displayed trace difference is nonzero.  The second inequality
holds for every finite set $\mathcal X$ of the specified actual
representations, with no two distinct members isomorphic, at every
place and Frobenius allowed in (SP.44).
\end{proposition}

\begin{proof}
Put $y=p^e$ and $x=p^{(e+1)/2}$.  Since $e\ge2$ and $p\ge5$,
$x\le y/2$ and $\sqrt p\,x\le y$, so (SP.43) is at most
$7y/2$.  For any three complex numbers of modulus one, the leading
term $p^{e+1}z_1$ has modulus $py$, whereas the sum of the trace and
the other three terms has modulus at most
\[
          \frac72y+y+p+1<py.
\]
The strict inequality follows from
$(p-9/2)p^e>p+1$, whose smallest case is $p=5,e=2$ and whose two
sides then separate monotonically.  Thus the difference is nonzero;
its modulus is less than $2py\le (py)^2=p^{2(e+1)}$.

For $u\in(\mathbb Z/k)^\times$ let
\[
 b(s,u)=\sum\beta_s-
          \bigl(p^{e+1}z_1^u+p^ez_2^u+pz_3^u+1\bigr),
 \qquad N_* =\prod_{s,u}b(s,u).
 \tag{SP.47}
\]
Every factor is an algebraic integer.  An automorphism of
$\mathbb C$ permutes the $s$'s and acts on all $k$-th roots by one
common unit modulo $k$; it therefore permutes the pairs $(s,u)$.
Consequently $N_*$ is fixed by every automorphism of its normal
number field over $\mathbb Q$, so is rational, hence an integer.
The nonvanishing just proved gives $N_*\ne0$.  Its logarithmic
size is at most
$2|\mathcal D|\varphi(k)(e+1)\log p$.
Use $e+2\le2(e+1)$ and $e+1\le M(e+1)$ to obtain (SP.45)--(SP.46).

For a fixed finite family $\mathcal X$, assign to $r$ its label
$s(r)$ and select the factor $b(s(r),1)$.  These selected factors
are distinct because equal labels imply isomorphism.  Equation
(SP.44) identifies their inverse $\iota$-images with exactly the
trace differences in (SP.46).  Every omitted factor is still an
algebraic integer and thus has $\ell$-adic norm at most one.
Taking the norm of the complete product proves the asserted
inequality.  The same argument with a singleton proves the
nonvanishing for every individual representation.  All the choices
defining $N_*$ preceded the auxiliary prime and the family, so no
uniformity has been exchanged in this argument.
\end{proof}

\subsection{The paired-count join and its remaining geometric input}

For precision, we state the outstanding input at the level at which
the preceding proofs can actually use it.  The following two
assertions remain to be reconstructed.
\begin{enumerate}
\item[(G1)] For every original host $h$, $h.\mathrm{Big}(t)$,
odd $\ell$, $0<e<t$ and $\psi\in h.X(c,e)$ in case II or III,
after adjoining the roots allowed in SP8 there is an actual
$D:\texttt{Raw}\,\Omega\,c$ with \texttt{Good} and
\texttt{OfCase}.  There are the source parameters $n',p_1,m_0$,
with $m_0\mid h.o$, $c_2\le n'$, $D.e=e$ and
$D.\psi=\operatorname{dirG}(\psi)$, for which
\[
 \begin{aligned}
 n_{
m II}&=\lceil\log_2\max
       (\ell h.N4^{n_0},\,2^{c_2+6}h.o\ell^3h.q^2)\rceil,\\
 n_{
m III}&=\lceil\log_2\max
       (\ell h.q\,p_{\!1}(\chi)4^{n_0},\,8\ell^2)\rceil,\\
 (n',p_1)&=(n_{
m II},1)\quad\hbox{or}\quad
                  (n_{
m III},p_{\!1}(\chi)),\\
 L&=\ell 2^{n'}h.Np_1,
 & k&=(\ell-1)2^{n'}(2h.o).
 \end{aligned}
 \tag{SP.48}
\]
Here $p_{\!1}(\chi)$ denotes the fixed source \texttt{pOne}, not
the variable $p_1$.  The characters $D.\psi,D.\xi,D.\mathrm{om}$
have $k$-th power one and satisfy the source \texttt{KillsMu}
conditions at level $L$; $D.S$ is the prime support of
$2\ell h.Np_1$.  With one common $\iota$, every actual irreducible
$D.\mathrm{EPt}$ representation satisfies the \texttt{AdmRAlt}
conditions used in SP9 at level $L$.
\item[(G2)] For this same $D$ there is an actual
$\mathcal M:\texttt{Model}\,D$ whose first-weight-zero algebra,
point, grading and equations satisfy (M1)--(M4) above.  Every
two-sided point of its equation scheme determines an irreducible
$D.\mathrm{EPt}$ representation.  The assignment is compatible
with the trace element $a_\sigma$:
\[
 x(a_\sigma)=0,\qquad
 \xi(a_\sigma)=\operatorname{tr}\rho_\xi(\sigma)
                       -\operatorname{tr}D.r_0(\sigma).
 \tag{SP.49}
\]
Isomorphic assigned representations come from the same
$\mathbb G_m$ orbit.  The trace element belongs to the degree-zero
algebra, and the valuation and coefficient map are those of the
original $D$, without a change of residual character or local
condition.
\end{enumerate}
These assertions explicitly include the construction of the
weighted model and its local-to-global comparison.  Naming
\texttt{LemmaP}, \texttt{CohomologyHyp} or an automorphy package
does not prove them.

\begin{theorem}[SP10: the complete arithmetic and coefficient-return join]
\label{claim:SP10}
Assume (G1), (G2), and the finite catalogue preceding SP9.
Then SP1--SP9, RV3 and RV1 imply the literal
\texttt{CountII}$\wedge$\texttt{CountIII}.  In particular, for
each primitive character $\chi$ there is one natural exponent
$C_\chi$, chosen before $h,t,\Omega,\ell,c,e,\psi$, such that
every original character in either case has finite strict Selmer
length and
\[
 \frac{\ell_O V(\Omega,c,e,\psi)}{v_O(\ell)}\log\ell
                         \le(\ell h.N e)^{C_\chi}.
 \tag{SP.50}
\]
This is a conditional endpoint.  The theorem closes the join after
the displayed geometric and catalogue inputs; it does not close
(G1), (G2), or the catalogue itself.
\end{theorem}

\begin{proof}
Choose $C,A$ from SP9, before any host or coefficient choices.
For the current $h,\ell,e,\psi$, put
$p=p_0(\ell h.Np_1)$, the least prime at least five not dividing
$\ell h.Np_1$.  Bertrand's theorem applied to $\ell h.Np_1+5$
gives
\[
                    5\le p\le2\ell h.Np_1+10.
 \tag{SP.51}
\]
In particular $p\nmid L$ and $p\notin D.S$.  Choose arithmetic
Frobenius $\sigma$ at $D$'s specified place over $p$.  Its
cyclotomic value is $p$.  Set
\[
 \begin{aligned}
 z_1&=\iota((D.\mathrm{om}^{-1}D.\xi D.\psi)(\sigma)),\quad
 z_2=\iota(D.\psi(\sigma)),\quad
 z_3=\iota((D.\mathrm{om}^{-1}D.\xi)(\sigma)).
 \end{aligned}
\]
They are $k$-th roots of unity, and direct substitution in the four
diagonal characters $D.r_0=(D.\mathrm{th},D.\Psi,D.\mathrm{al},1)$
shows that its trace is
$\iota^{-1}(p^{e+1}z_1+p^ez_2+pz_3+1)$.
SP9 thus makes every irreducible trace difference nonzero.
By (SP.49), $a_\sigma$ satisfies (M5).

Apply SP6 to obtain a finite cover $T$ of two-sided points and the
cotangent inequality.  Remove redundant representatives until
there is one per orbit.  Its assigned representations form a
finite set $\mathcal X$ with no two distinct members isomorphic:
this is exactly the last orbit assertion of (G2).  The value of
$a_\sigma$ is constant on each orbit.  SP7 and SP6 now give
\[
 \ell_O(D.\mathrm{Sha})
       \le\ell_OC_{-1}(F)
       \le2\sum_{r\in\mathcal X}
                 v_O(\operatorname{tr}r(\sigma)
                                  -\operatorname{tr}D.r_0(\sigma)).
 \tag{SP.52}
\]
For an algebraic $\ell$-adic number $z\ne0$ the extension of the
normalized $O$-valuation is
$v_O(z)=-v_O(\ell)\log|z|_\ell/\log\ell$.
Consequently (SP.52) is the source's actual $D.\mathrm{Ineq}(\sigma)$,
with its exact factor two and its finite-length assertion.
The product estimate (SP.46) then yields
\[
 \frac{\ell_O(D.\mathrm{Sha})}{v_O(\ell)}\log\ell
 \le -2\log|N_*|_\ell
 \le 2\log|N_*|
 \le 2Ck[L(e+1)]^A\log p.
 \tag{SP.53}
\]
The middle inequality is just
$v_\ell(N_*)\log\ell\le\log|N_*|$ for a nonzero integer;
the last also uses $\varphi(k)\le k$.  This explains why both the
integrality and the order of the quantifiers in SP9 are needed.

The original Selmer group injects into $D.\mathrm{Sha}$ even
though $D.S$ can contain additional primes.  Indeed, at a prime
$p\nmid\ell h.N$ the original character is unramified, and its
negative-twist Frobenius value is $\zeta p^{-e}$ for a root of
unity $\zeta$.  It is not one since $e>0$.  On $E/O$ multiplication
by $\zeta p^{-e}-1\ne0$ is surjective.  Thus the unramified
cohomology
$(E/O)/(\zeta p^{-e}-1)(E/O)$ vanishes, and restriction from
decomposition-group cohomology to inertia is injective.
Every original class already vanishes on inertia at every place;
it therefore vanishes on each added decomposition group as well.
At a prime dividing $\ell h.N$ this was already its defining
strict local condition.  This proves the exact inequality called
\texttt{shaEq}, rather than an unwarranted equality, on the same
$\Omega$ and all its chosen valuations.

It remains to verify a single exponent uniformly in the host.
Write $o=\operatorname{ord}\chi$, $P=p_{\!1}(\chi)\ge1$,
$f=\varphi(4o)$, and $X=\ell h.Ne$.  The host's numerical and
character conditions give
\[
 \begin{gathered}
 X\ge16,\quad \ell,h.N,e,h.o,h.q\le X,\quad
 2^{c_2}\le h.N,\quad 2^{n_0}<\ell^f. 
 \end{gathered}
 \tag{SP.54}
\]
For clarity, the last bound uses no analytic theorem.  The order
of $\theta_1$ is prime to $\ell$ and divides $2o$.
Here is the finite-group reason for the latter statement.
Under the Chinese remainder decomposition, a unit supported on an
odd prime-power factor dividing $h.q$ has its $o$-th power in
$h.H$: it is trivial at the other primary factors and its image
under $\chi$ has order dividing $o$.  A unit supported on any
other odd prime-power factor has its square in $h.H$, by the
two square-subgroup conditions defining $h.H$.  Since
$h.H\subset\ker\psi$, every odd primary component of $\psi$
has order dividing $2o$.  Taking its prime-to-$\ell$ part preserves
that divisibility, and removing its $2$-primary-modulus component
leaves exactly $\theta_1$.  The product of the remaining primary
components therefore has order dividing $2o$.
It follows that
$\operatorname{lcm}(4,\operatorname{ord}\theta_1)$ divides $4o$;
the order $f_0$ of $\ell$ modulo this modulus divides its unit
group order, hence is at most $\varphi(4o)$.  Therefore
$2^{v_2(\ell^{f_0}-1)}\le\ell^{f_0}-1<\ell^f$.
The other clauses use $16\mid h.N$, $h.q\mid h.N$,
$h.o=o\le h.q$, and the conductor divisibility
$2^{c_2}\mid h.N$, all on the original host and character.

The inequality $2^{\lceil\log_2 m\rceil}\le2m$ for positive
integers $m$ gives, from (SP.48) and (SP.54),
\[
 2^{n_{
m II}}\le X^{2f+10},\qquad
 2^{n_{
m III}}\le X^{2f+P+3}.
 \tag{SP.55}
\]
For the first, the two arguments of the maximum are at most
$X^{2f+2}$ and $64X^7\le X^9$, and the extra factor two is at
most $X$.  For the second they are at most $X^{2f+P+2}$ and
$8X^2\le X^3$, using $P\le X^P$.  Set $a=2f+P+13$; either
choice of $n'$ then satisfies $2^{n'}\le X^a$.  Directly,
\[
 \begin{aligned}
 k&\le X^{a+3},& L(e+1)&\le X^{a+P+4},\\
 2C&\le X^{2C},& \log p&\le p\le X^{P+3}.
 \end{aligned}
 \tag{SP.56}
\]
For the last bound use (SP.51), $\ell h.Np_1\le X^{P+2}$
and $X\ge16$.  Hence one explicit, deliberately generous choice is
\[
             C_\chi=2C+a+3+A(a+P+4)+P+3.
 \tag{SP.57}
\]
It depends only on $\chi,C,A$.  Combining (SP.53), the original
Selmer inclusion, and (SP.56) proves (SP.50) over the enlarged
coefficient ring.  SP8 returns it to the original $O$ with the
same exponent, case, order, kernel, conductor, parity and
normalized length.  The same $C_\chi$ works for the union of the
two cases, so gives the two literal count predicates separately.
\end{proof}

\paragraph{Precisely what is still open in the shared count branch.}
The algebraic leaves \texttt{LemmaM'} and \texttt{LemmaM} have been
proved in SP1 and SP6.  SP7 also proves \texttt{shaCoc} and
\texttt{mPrime} for the actual model; RV7 supplies both
\texttt{SoV} and \texttt{SoFin}, while BR7 supplies the literal
\texttt{CorX}.  These remove complete suppliers from the source's
\texttt{sp4Main\_alt} input list.  The following are still
substantive inputs to the human reconstruction:
\begin{enumerate}
\item \texttt{SymplecticFourImageDense}, including both the finite
weight-zero algebra and exclusion of the endoscopic locus, with
one embedding valid for every actual irreducible point.
\item The entire \texttt{CohomologyHyp} package.  LK2--LK4 and LD3
cover identified finite local and finite duality pieces; they do
not by themselves identify the arbitrary period-ring records in
its Bloch--Kato local clause, or prove its remaining global
finiteness, Euler, orthogonality and Greenberg--Wiles clauses.
\item The application of the literal \texttt{LemmaP} to the actual
Galois data, with its framed weighted representation functor,
specified local conditions, relation weights and exact
cohomological dimension counts.  The universal algebraic
implication \texttt{LemmaP} is handled separately by the framed
deformation component; this chapter does not infer its actual
hypotheses or the resulting \texttt{exModel} comparison from that
implication alone.  The cumulative ledger records the status and
supplier of that separate algebraic theorem.
\item The finite catalogue used in SP9, including its polynomial
size, algebraicity, conjugation, purity and classification
properties for the literal \texttt{AdmRAlt}.  The source derives
it from a larger automorphic and local-level package; its name
\texttt{Hypotheses'} is not a classical theorem citation.
\item \texttt{RigiditySymplecticFour} at these actual points.  Its
source producer uses the specified unramified deformation inputs
and adjoint lifting theorem; no generic rigidity assertion has
been substituted.
\item The internal construction and comparison edges in (G1)--(G2):
\texttt{dataEx}, \texttt{exModel}, the verification of (M1)--(M4),
the two-sided point and conjugacy-orbit comparison, and the
\texttt{level}/\texttt{safe4'} passage to the exact admissible
representation class.  Their required inputs and source ranges
are recorded separately in the provenance ledger.  They remain
open human proof obligations even when their larger source input
packages have been named.
\end{enumerate}
The source's exact top-level implication is recorded for navigation;
SP10 is the readable conditional join actually proved here.  Thus
this chapter does not assert that \texttt{CountII} or
\texttt{CountIII}, the original-family \texttt{PointHV}, or RH/GRH
has now been proved without the remaining suppliers.

\paragraph{Classical reference for the prime choice.}
The only additional number-theoretic theorem in the last join is
Bertrand's theorem, in the form: for every integer $n>1$ a prime
lies strictly between $n$ and $2n$.  See P.~Erd\H{o}s,
\emph{Beweis eines Satzes von Tschebyschef}, Acta Sci. Math.
(Szeged) \textbf{5} (1932), 194--198,
\url{https://users.renyi.hu/~p_erdos/1932-01.pdf}.
It is applied only with $n=\ell h.Np_1+5$ to obtain (SP.51).

